\documentclass[openany]{ctexbook}
\usepackage{amsmath,amssymb}
\usepackage[hidelinks]{hyperref}
\newcommand{\note}[1]{\footnote{译注：#1}}
\title{代数几何基础第五卷预备笔记\\9.1至9.2：更换射影嵌入}
\author{据用户上传英译编辑稿译出}
\date{2026年11月15日}
\begin{document}
\maketitle
\chapter*{范围与校订状态}
本批连续翻译底本第36页第9节标题、9.1、命题9.2及其后附记，保留原脚注52。下一处为第37页命题9.3标题及陈述。本批为独立增补，未并入主合订本。9.2所引8.8的缺失条件未恢复，其附记也表明原陈述仍待调整；本译文不将该命题作为已经独立审定的结论。前批维数、对偶及锥性措辞等疑点均未因此消除。
\chapter*{§9\quad 更换射影嵌入}
\noindent\textbf{9.1.}
对每个整数 $n>0$，令
\[
P(n)=\mathbb P(\operatorname{Sym}^{n}(E)).
\]
有一个显然的浸入 $u_n:P\to P(n)$，因为在 $S$ 的每个仿射开集上，$\mathcal O_P(n)$ 都由其截面生成，并且
\[
p_*\mathcal O_P(n)\simeq\operatorname{Sym}^{n}(E),
\]
其中 $p:P\to S$ 是投影。\footnote{原脚注52：“无法辨认。”底本此处印为 $p_*(0_p(n))\operatorname{Sym}^{n}(E)$，缺关系符，前式还将 $E$ 印作小写 $e$。译校说明：正文统一记号，并依射影丛的商约定写出标准典范同构；这是数学上可核实的整理，不声称恢复了手稿原字迹。由评估满射 $p^*\operatorname{Sym}^{n}(E)\twoheadrightarrow\mathcal O_P(n)$ 得到相对 Veronese 浸入；在 $E$ 平凡化的开集上，其坐标为全部 $n$ 次单项式。}
若 $f:X\to P$ 是非分歧态射（分别地，是浸入），则 $u_n f:X\to P(n)$ 也如此。有时，研究 $X$ 时用 $u_n f$ 代替 $f$ 是有利的，可以避免 $f$ 在某些方面所具有的十分特殊、且有时令人棘手的行为。
（这种特殊情形的一个例子是8.12(b)中所指出的〔底本原有“sic”标记〕情形：$Y^{\mathrm{sing}}\to P^\vee$ 的像具有维数 $r-1$，但产生的域扩张不可分。另一个例子由 $\mathbb P^3$ 中的二次曲面给出：其所有奇异超平面截面都几何可约，尽管 $X$ 本身几何不可约。）\note{这里按原文保留二次曲面这一例子的概括措辞，未扩大为任意特征、任意退化二次曲面的无条件断言。前文几何普通奇点的术语约定沿用原稿；不将“几何可约”与“非约化”混同。}

\medskip
\noindent\textbf{命题9.2〔原稿陈述，所引条件待核〕.}
设 $S=\operatorname{Spec}(k)$，$X$ 在 $k$ 上光滑，且 $f:X\to P$ 非分歧。取 $n\ge2$，并考虑 $f_n=u_n f$。原稿断言：$f_n:X\to P(n)$ 满足8.8中的等价条件；特别地，在超平面参数空间 $P(n)^\vee$ 中，除去一个余维至少为2的集合后，对于其中的点 $\xi$，相应的超平面截面 $Y_\xi$ 或者光滑，或者只有有限个非光滑点，且这些点都是几何普通奇点。若 $f$ 是浸入，则这样的奇点至多只有一个，并且它在 $k(\xi)$ 上有理。\note{底本写 $\xi\in P(n)$。由于 $\xi$ 参数化 $P(n)$ 中的超平面，正文订为 $P(n)^\vee$，与9.1所用 $P^\vee$ 及全篇的超平面参数约定一致。此为参数空间记号订正，并非补全8.8条件。底本此处未给出完整证明；下文仅有关于8.8(ii bis)的核验提示。}

\medskip
\noindent\textbf{附记（N.B.）.}
应当以一种不排除 $f$ 并非浸入这一情形的方式来陈述8.8。在8.8的条件(ii bis)下，核验本质上是平凡的。\note{“核验本质上是平凡的”为原稿评语，不是本次独立审校的结论。现有底本缺失8.8及其(ii bis)条件，故无法仅凭这句完成核验，也不能据此声称已恢复缺文。}
无疑，应当在9.1中明确说明：$X$ 相对于 $f_n$ 的超平面截面，无非是用次数为 $n$ 的超曲面代替超平面所得的 $X$ 的“截面”。
\end{document}
